\section{When Is a Public Hub Efficient?}
\label{sec:hub_efficiency}

This section characterizes the conditions under which a public hub is socially efficient \emph{as a stand-alone local policy}. Thus, throughout this section, we set
\[
b=0.
\]
The purpose is to isolate the hub's value through local coordination inside region $A$ before turning in Section 6 to the case in which the hub is combined with a cross-regional pipeline policy.

\subsection{A Benchmark Efficiency Criterion}

From \eqref{eq:DeltaH}, the net welfare gain from establishing a hub without pipeline policy is
\begin{equation}
\Delta_H
=
v_L P_A \bigl(m_A^1-m_A^0\bigr)
\Bigl[(1+\beta)-\beta\delta\bigl(m_A^1+m_A^0\bigr)\Bigr]
-F.
\label{eq:DeltaH_repeat}
\end{equation}

It is useful to define the hub's \emph{dynamic viability factor} as
\begin{equation}
\Lambda_A
:=
(1+\beta)-\beta\delta\bigl(m_A^1+m_A^0\bigr).
\label{eq:Lambda_def}
\end{equation}
Then \eqref{eq:DeltaH_repeat} becomes
\begin{equation}
\Delta_H
=
v_L P_A \bigl(m_A^1-m_A^0\bigr)\Lambda_A - F.
\label{eq:DeltaH_Lambda}
\end{equation}

The term $m_A^1-m_A^0$ captures the hub's immediate local coordination value: how much it raises the effective matching rate within region $A$. By contrast, $\Lambda_A$ captures whether those short-run gains remain socially valuable once second-period local redundancy is taken into account. A larger $\delta$ lowers $\Lambda_A$, making the hub less attractive from a dynamic welfare perspective.

The hub is socially efficient if and only if
\[
\Delta_H \ge 0.
\]

\begin{proposition}
\label{prop:hub_efficiency}
Consider the benchmark case without pipeline policy, $b=0$.

\begin{enumerate}[leftmargin=2em]
    \item If
    \[
    \Lambda_A \le 0,
    \]
    then the hub is never socially efficient:
    \[
    \Delta_H < 0
    \qquad
    \text{for all } N_A \ge 2 \text{ and } F>0.
    \]

    \item If
    \[
    \Lambda_A > 0,
    \]
    then the hub is socially efficient if and only if
    \begin{equation}
    P_A
    \ge
    \frac{F}{v_L\bigl(m_A^1-m_A^0\bigr)\Lambda_A}.
    \label{eq:pair_threshold}
    \end{equation}
    Equivalently, treating $N_A$ as continuous, the hub is efficient if and only if
    \begin{equation}
    N_A \ge N_H^*
    :=
    \frac{1}{2}
    \left(
    1+
    \sqrt{
    1+
    \frac{8F}{v_L\bigl(m_A^1-m_A^0\bigr)\Lambda_A}
    }
    \right).
    \label{eq:Nstar}
    \end{equation}
    If $N_A$ is restricted to integers, replace $N_H^*$ by its ceiling.
\end{enumerate}
\end{proposition}

Proposition \ref{prop:hub_efficiency} has a simple interpretation. A hub creates value only if the gross local coordination benefit,
\[
v_L P_A \bigl(m_A^1-m_A^0\bigr),
\]
survives the dynamic penalty from repeated local interaction, summarized by $\Lambda_A$, and is large enough to cover the fixed cost $F$. Thus, a hub may fail to be efficient for two conceptually distinct reasons. First, the target region may be too thin, so that $P_A$ is too small. Second, even if the region is not especially thin, the redundancy parameter $\delta$ may be large enough to make repeated local interaction dynamically unproductive.

A useful implication of the two-region formulation is that the stand-alone criterion in \eqref{eq:DeltaH_Lambda} depends only on \emph{local} primitives. The size of the outside pool $N_B$, the baseline outside access $q_0$, and the surplus attached to cross-regional links $v_X$ affect welfare levels in \eqref{eq:W00} and \eqref{eq:W10}, but they cancel out of the incremental comparison between a hub alone and no hub.

\subsection{Equivalent Threshold Characterizations}

Proposition \ref{prop:hub_efficiency} can be rewritten in several equivalent ways, which are useful for interpretation.

First, for given $(N_A,m_A^0,m_A^1,\beta,\delta)$, define the critical fixed cost
\begin{equation}
F_H^*(N_A)
:=
v_L P_A \bigl(m_A^1-m_A^0\bigr)\Lambda_A.
\label{eq:Fstar}
\end{equation}
Then the hub is efficient if and only if
\[
F \le F_H^*(N_A).
\]

Second, for given $(N_A,m_A^0,m_A^1,\beta,F)$, define the critical redundancy rate
\begin{equation}
\delta_H^*(N_A)
:=
\frac{1+\beta}{\beta\bigl(m_A^1+m_A^0\bigr)}
-
\frac{F}{\beta v_L P_A\bigl(m_A^1-m_A^0\bigr)\bigl(m_A^1+m_A^0\bigr)}.
\label{eq:deltastar}
\end{equation}
Then the hub is efficient if and only if
\[
\delta \le \delta_H^*(N_A),
\]
provided the right-hand side is positive.

These equivalent thresholds make clear that the hub's stand-alone efficiency is governed by five local primitives: the size of the target ecosystem $N_A$, the hub's matching improvement $m_A^1-m_A^0$, the surplus per realized local interaction $v_L$, the fixed cost $F$, and the speed of local redundancy $\delta$.

\subsection{Comparative Statics}

The next proposition summarizes how the hub's net welfare contribution responds to the primitive parameters.

\begin{proposition}
\label{prop:comparative_statics_hub}
Suppose $\Lambda_A>0$. Then the hub-alone welfare gain $\Delta_H$ satisfies
\begin{align}
\frac{\partial \Delta_H}{\partial N_A}
&=
\frac{v_L}{2}(2N_A-1)\bigl(m_A^1-m_A^0\bigr)\Lambda_A
>0,
\label{eq:dDelta_dN}
\\[0.4em]
\frac{\partial \Delta_H}{\partial F}
&=
-1
<0,
\label{eq:dDelta_dF}
\\[0.4em]
\frac{\partial \Delta_H}{\partial \delta}
&=
-\beta v_L P_A \bigl(m_A^1-m_A^0\bigr)\bigl(m_A^1+m_A^0\bigr)
<0.
\label{eq:dDelta_ddelta}
\end{align}
Moreover, under the regularity condition \eqref{eq:regularity},
\begin{align}
\frac{\partial \Delta_H}{\partial m_A^1}
&=
v_L P_A\Bigl[(1+\beta)-2\beta\delta m_A^1\Bigr]
>0,
\label{eq:dDelta_dm1}
\\[0.4em]
\frac{\partial \Delta_H}{\partial m_A^0}
&=
- v_L P_A\Bigl[(1+\beta)-2\beta\delta m_A^0\Bigr]
<0.
\label{eq:dDelta_dm0}
\end{align}
\end{proposition}

Proposition \ref{prop:comparative_statics_hub} formalizes several intuitive points.

First, the hub is more likely to be efficient in a thicker target ecosystem. Since
\[
P_A=\frac{N_A(N_A-1)}{2},
\]
the number of potential local collaborative pairs rises more than proportionally with $N_A$. This is the sense in which thin regional ecosystems face a structural disadvantage: even if the hub improves local matching, there may simply be too few potentially valuable local interactions for the gain to cover the fixed cost.

Second, faster local redundancy makes the hub less attractive. This is the model's formal representation of local lock-in or repeated interaction among the same actors inside region $A$. A hub is therefore not justified by short-run coordination gains alone; its dynamic viability depends on whether those gains persist.

Third, a stronger local matching effect of the hub, measured by a larger $m_A^1$, raises the hub's welfare contribution, whereas a higher baseline local matching rate $m_A^0$ lowers its incremental value. Thus, a public hub is most useful where pre-existing local coordination is weak but where the hub can substantially lower interaction frictions.

\subsection{Discussion}

This section yields two implications that will be important later.

The first is that hub efficiency is not determined by local scale alone. Even a relatively large target region may not justify a hub if local redundancy builds quickly. Conversely, even a thin region may justify a hub if the coordination gain is large enough and dynamic redundancy is limited.

The second is that hub efficiency depends on a dynamic object, $\Lambda_A$, rather than only on the static matching improvement $m_A^1-m_A^0$. This foreshadows the central message of the paper: a public hub can look attractive on short-run local coordination grounds and yet be dynamically inefficient once the risk of repeated interaction within the same regional pool is taken seriously.

For that reason, the next section distinguishes explicitly between short-run success and total dynamic efficiency.
